% =================================================================
\section{Exam technique and the errors that cost marks}
\label{sec:technique}
% =================================================================

\begin{tip}[What examiners keep saying (themes from VCAA examination reports)]
\begin{itemize}
  \item \textbf{Show the key step.} In Exam 1, and in Exam 2 questions worth more than 1 mark, an answer alone often earns only part of the marks. Write the equation you are solving, the derivative or the integral with its bounds.
  \item \textbf{Use correct notation.} $\dydx=\dots$ or $f'(x)=\dots$, not ``$=$'' chains joining unequal things; integrals need $dx$; indefinite integrals need $+c$; write $\loge$, not calculator syntax; use set and interval notation correctly.
  \item \textbf{Answer the actual question.} Give coordinates when asked for a point, the \emph{value} when asked for the maximum value (not the $x$ where it occurs), all values of a parameter when asked for ``values'', units, and a sentence in context.
  \item \textbf{Exact means exact.} When decimals are asked for, round only at the end and to the stated accuracy.
  \item \textbf{``Hence''} means use the previous part. ``Show that'' means every step to the given answer, with no gaps.
\end{itemize}
\end{tip}

\subsection{Topic-by-topic errors to avoid}

\begin{warn}[Functions and graphs]
\begin{multicols}{2}
\begin{itemize}
  \item Missing asymptote equations, open/closed endpoints, or coordinates on sketches.
  \item Graphs that bend the wrong way near asymptotes or cross an asymptote they can't (e.g.\ $y=e^x$ crossing $y=0$).
  \item Dilation factor: $f(2x)$ is a factor of $\frac12$ from the $y$-axis. Translation direction: $f(x+3)$ moves \emph{left}.
  \item Doing transformations in the wrong order, or not factorising inside the brackets before reading off.
  \item Writing $f^{-1}$ for $\frac1f$; forgetting the domain of $f^{-1}$ (it is the range of $f$).
  \item Composite domain: $d_{f\circ g}=d_g$, not the implied domain of the simplified rule.
  \item CAS in degree mode.
\end{itemize}
\end{multicols}
\end{warn}
\begin{warn}[Algebra and equations]
\begin{multicols}{2}
\begin{itemize}
  \item Missing solutions in circular equations because the interval for $ax+b$ wasn't adjusted.
  \item Dividing by an expression that might be zero (e.g.\ $\sin x$) and losing solutions. Factorise instead.
  \item Not rejecting invalid solutions of log equations (check that every argument is $>0$).
  \item $\loge(a+b)\ne\loge a+\loge b$; \ $e^{x+y}=e^xe^y$.
  \item Simultaneous equations: when the determinant is $0$, you still have to check whether there are no solutions or infinitely many.
  \item The $k=0$ case in ``$kx^2+\dots$'' (it is not a quadratic then).
\end{itemize}
\end{multicols}
\end{warn}
\begin{warn}[Calculus]
\begin{multicols}{2}
\begin{itemize}
  \item Forgetting the inner derivative in the chain rule ($\ddx e^{3x}=3e^{3x}$).
  \item Quotient rule in the wrong order: it is $\frac{vu'-uv'}{v^2}$.
  \item $\int\frac{1}{ax+b}\dx=\frac1a\loge(ax+b)+c$ ($ax+b>0$): don't drop the $\frac1a$. Likewise $\int\sin(ax)\dx=-\frac1a\cos(ax)$: watch the sign.
  \item Integrating across an $x$-intercept and calling the signed integral the ``area''.
  \item Area between curves: use top minus bottom, and split where the curves cross.
  \item Confusing the average value $\frac{1}{b-a}\int_a^bf$ with the average rate of change $\frac{f(b)-f(a)}{b-a}$.
  \item Not justifying the nature of a stationary point; ignoring endpoints for absolute max/min.
  \item Strictly increasing intervals that join across an asymptote or point of discontinuity.
  \item Trapezium rule: wrong strip width $h$, or not doubling the interior values. Newton's method: rounding $x_n$ too early.
\end{itemize}
\end{multicols}
\end{warn}
\begin{warn}[Probability and statistics]
\begin{multicols}{2}
\begin{itemize}
  \item $\Nd(\mu,\sigma^2)$: the second parameter is the \emph{variance}. Enter $\sigma$ (not $\sigma^2$) into CAS.
  \item Discrete inequalities: $\Pr(X>2)=1-\Pr(X\le2)$, not $1-\Pr(X\le3)$.
  \item Conditional probability: the denominator is the probability of the \emph{given} event.
  \item Using the binomial distribution when the trials aren't independent (sampling without replacement), or not stating the binomial conditions when asked to justify.
  \item pdf questions: checking only one condition; mixing up the median (area $0.5$) and the mean ($\int xf(x)\dx$).
  \item Confidence intervals: incorrect interpretation; rounding $n$ \emph{down} (always round up); using the normal approximation when the question says to use the binomial (or vice versa).
\end{itemize}
\end{multicols}
\end{warn}

\subsection{Typical question types}

\begin{cols}
\begin{keyres}{Exam 1 (no technology): practise until automatic}
\begin{itemize}
  \item Differentiate (product, quotient, chain) and evaluate at a point.
  \item ``Differentiate \dots\ hence find an anti-derivative of \dots''
  \item Anti-derivatives of $f(ax+b)$; definite integrals with exact values; area between curves.
  \item Solve a circular, exponential or log equation exactly over an interval.
  \item Inverse and composite functions with domains and ranges.
  \item Sketch a transformed exponential, log, power or polynomial graph with features labelled.
  \item Tangent and normal lines; stationary points; a simple optimisation.
  \item Discrete table with an unknown, mean and variance; binomial with small $n$; pdf: find $k$, a probability, a median; normal symmetry using given probabilities; small-sample $\Phat$; CI with $z\approx2$.
  \item Simultaneous equations with a parameter; one or two steps of Newton's method; the trapezium rule; trace pseudocode.
\end{itemize}
\end{keyres}
\nextcol
\begin{keyres}{Exam 2 (CAS): what Section B usually looks like}
\begin{itemize}
  \item \textbf{Function and calculus modelling:} a curve, track, tank or population model. Find features, rates and optimal values; areas; join curves smoothly (hybrid functions).
  \item \textbf{Circular modelling:} tides, temperature or a Ferris wheel. Find amplitude, period and translations; solve equations for times; average values.
  \item \textbf{Probability:} discrete $\to$ binomial $\to$ continuous pdf $\to$ normal $\to$ sample proportions and CIs, often linked through one context, with conditional probability.
  \item \textbf{Families and parameters:} a function with a parameter $a$ or $k$. Find the values for which there are $n$ solutions, a tangent, an inverse, etc.
  \item \textbf{Transformations and inverses} combined with calculus (e.g.\ the area between $f$ and $f^{-1}$).
\end{itemize}
\textbf{Section A (MCQ):} quick checks by graphing, substituting the options, or testing special values. About 1.5 minutes each.
\end{keyres}
\end{cols}

\begin{tip}[Final-week routine]
\begin{enumerate}
  \item Do full past exams under timed conditions. VCAA papers from 2023 onwards match this study design. Older papers (2016--2022) are still good practice for most topics: check any unfamiliar question against the checklist in Section~\ref{sec:checklist}.
  \item Read the matching VCAA examination report for each paper. It lists exactly where marks were lost.
  \item Re-do every question you got wrong 2--3 days later, without looking.
  \item For Exam 2, finalise your bound reference: this sheet, CAS syntax notes and your own worked solutions. Tab it by topic so you can find things in seconds.
  \item Use the checklist in the next section to find your last gaps.
\end{enumerate}
\end{tip}
